\section{Introduction}
\label{sec:intro}

The modern paradigm of deploying deep neural networks at the edge is governed by a strict, multi-faceted set of systems constraints: latency must be sub-millisecond, memory footprints must fit within small embedded SRAMs, and power budgets are capped at milliwatts. To achieve these goals, machine learning practitioners utilize model compression and low-precision integer quantization (typically 8-bit integer [INT8] for activations and 4-bit integer [INT4] for weights) as standard industry practices \cite{jacob2018quantization, gholami2021survey}. These techniques yield massive speedups, reduce memory bandwidth, and decrease silicon cost, making high-performance models viable for real-world devices such as mobile phones, automotive microcontrollers, and edge gateways.

Concurrently, there is a growing need for edge systems to handle highly non-stationary, multi-task streaming queries. Rather than deploying multiple heavy, separate models for each downstream task, dynamic model ensembling and weight-merging in latent coordinate spaces have gained significant traction \cite{wortsman2022model, ainsworth2022git, choshen2022fusing}. Built upon the high-fidelity Coordinate Sandbox (ICS) environment, methods like Stateless Activation-Based Latent Ensembling (\textsc{Sable}) \cite{sable2024stateless}, continuous-time chemical kinetics routing (\textsc{ChemMerge}) \cite{chemmerge2025continuous}, and simple exponential moving averages (\textsc{Momentum-Merge}) \cite{momentummerge2025exponential} demonstrate that task-specific expert adapters can be dynamically blended on the fly using intermediate activation cues. Under standard 32-bit floating-point (Float32) environments, these methods achieve remarkable continuous adaptation with zero training-time overhead, establishing a flexible and robust performance ceiling.

\begin{figure}[t]
  \centering
  \includegraphics[width=\linewidth]{comparison_plot.png}
  \caption{\textbf{Quantization Collapse and Recovery under Varying Entanglement.} At large sample calibration ($N_{\text{cal}} = 4000$), standard ensembling baselines (SABLE, ChemMerge, Momentum-Merge) achieve substantial classification gains over static Uniform Merging in full-precision (Float32). However, under naive INT8/INT4 quantization, their performance collapses directly to the uniform baseline (SABLE Quantized-Naive). Our proposed \textbf{QA-Merge} successfully recovers this performance, preserving high ensembling accuracy across all entanglement levels $\rho \in [0.0, 0.5]$.}
  \label{fig:quantization_collapse_intro}
\end{figure}

Despite their algorithmic elegance, a major chasm remains: \emph{these ensembling systems are built under the assumption of high-precision floating-point coordinate representations.} When these algorithms are compiled and deployed under extreme low-precision limits (INT8 activations and INT4 blending weights), we observe a catastrophic failure mode which we term the \textbf{Quantization Collapse}. 

As illustrated in Figure~\ref{fig:quantization_collapse_intro}, standard, naive quantization completely erases the benefits of dynamic ensembling. SABLE and ChemMerge drop directly to the accuracy level of simple static Uniform Merging. This collapse occurs because the rounding operators inherent to uniform quantization act as aggressive, non-differentiable step filters over intermediate activation coordinates. This rounding noise destroys the subtle representation boundaries that routing networks rely on, leading to overlapping task centroids, dead gradients, and highly biased, near-one-hot misrouting trajectories.

From a pragmatist's perspective, dynamic model ensembling is only useful if it can be reliably deployed under standard production constraints. Building on this core philosophy, we present \textbf{QA-Merge} (\textbf{Q}uantization-Aware \textbf{Merge}), a hardware-compatible suite of techniques designed specifically to stabilize dynamic activation-space ensembling within low-precision pipelines. Rather than introducing complex physical or high-order theoretical abstractions, QA-Merge introduces four practical, computationally light mechanisms that directly address the root causes of representation collapse:

\begin{itemize}
    \item \textbf{Quantized Centroid Calibration (QCC):} Calibrates task-specific centroids directly in the target quantized integer representation space during offline few-shot calibration, maximizing task separation and directionality.
    \item \textbf{Straight-Through Estimator (STE) Gating:} Employs the Straight-Through Estimator \cite{bengio2013estimating} to bypass the non-differentiable rounding operator during calibration, enabling robust, end-to-end backpropagation over discrete boundaries.
    \item \textbf{Error-Feedback Trajectory Stabilization (EF-Smooth):} Tracks blending coefficient rounding errors layer-by-layer and injects them back as a high-pass feedback correction, stabilizing representation trajectories under strict 4-bit constraints.
    \item \textbf{Activation Error Feedback (AEF):} Accumulates representation rounding errors layer-by-layer and feeds them back residually, completely overcoming the Small-Step Quantization Bottleneck on coarse INT8 activation grids.
\end{itemize}

We evaluate QA-Merge inside the 14-layer Coordinate Sandbox (ICS) environment under a series of non-stationary, multi-task streaming scenarios. Our empirical findings demonstrate that QA-Merge successfully recovers full-precision ensembling gains under extreme INT8/INT4 constraints across both small-sample ($N_{\text{cal}}=64$) and large-sample ($N_{\text{cal}}=4000$) regimes. This performance recovery is achieved with zero extra parameter capacity and robust trajectory stability. 

In summary, the primary contributions of this work are:
\begin{enumerate}
    \item We identify and empirically deconstruct the \emph{Quantization Collapse} phenomenon in dynamic coordinate-space model ensembling, illustrating how discrete rounding noise collapses high-precision routing manifolds.
    \item We formalize \emph{QA-Merge}, introducing QCC, STE Gating, EF-Smooth, and AEF as a unified, hardware-friendly framework to stabilize representations and trajectories under low-precision constraints.
    \item We conduct exhaustive evaluations in the Coordinate Sandbox across standard visual signatures, proving that QA-Merge successfully closes the gap to the Float32 ensembling ceiling while preserving low trajectory jitter.
\end{enumerate}

The remainder of this paper is organized as follows. Section~\ref{sec:related} reviews the related literature on model merging, quantization, and edge serving. Section~\ref{sec:method} details the mathematical formulation and architectural design of QA-Merge. Section~\ref{sec:experiments} presents our comprehensive empirical results, and Section~\ref{sec:conclusion} concludes with a discussion of pragmatic future directions.
